\documentclass[10pt, a4paper]{article}
\usepackage[utf8]{inputenc}
\usepackage[T1]{fontenc}
\usepackage{amsmath, amssymb}
\usepackage{geometry}
\usepackage{xcolor}

% Configuration de la page
\geometry{
 a4paper,
 total={170mm,257mm},
 left=20mm,
 top=20mm,
}

% Commandes personnalisées pour l'oral
\newcommand{\IdeeP}[1]{\noindent\textbf{\textcolor{blue}{IDÉE PRINCIPALE : #1}}\par}
\newcommand{\ChiffreCle}[1]{\noindent\textbf{\textcolor{red}{\textbullet\ Chiffre Clé : #1}}\par}
\newcommand{\SousAxe}[1]{\vspace{0.5em}\noindent\underline{\textbf{#1}}\par}
\newcommand{\Transition}[1]{\vspace{1em}\noindent\textit{\textcolor{gray}{[Transition : #1]}}\par}

\title{\textbf{Prise de Notes Oral : Accès aux Soins en France}}
\author{MIGNARD Maël}
\date{\today}

\begin{document}

\maketitle
\thispagestyle{empty} % Supprime le numéro de page sur la première page

% ----------------------------------------------------------------------
% SECTION PRINCIPALE : L'ACCÈS AUX SOINS EN FRANCE
% ----------------------------------------------------------------------

\section*{L'Accès aux Soins en France : Un Système à Quatre Niveaux de Fracture}

\IdeeP{Le système français garantit l'accès aux droits, mais la réalité socio-économique et territoriale crée de fortes inégalités d'accès.}

\subsection*{\underline{1. La Carte Vitale et le Tiers Payant}}

\IdeeP{La Carte Vitale est le pilier de la simplification, essentiel pour les plus modestes.}
\SousAxe{Fonction principale : Facilitation administrative}
\begin{itemize}
    \item **Télétransmission :** Déclenchement rapide du remboursement.
    \item **Tiers Payant :** Dispense d'avance de frais pour la part Sécurité Sociale (AMO).
    \item **CSS :** Vital pour les bénéficiaires de la Complémentaire Santé Solidaire, annulant la barrière financière immédiate.
\end{itemize}
\ChiffreCle{Principe fondamental : L'assurance maladie publique est un puissant outil de redistribution en nature.}

\Transition{Malgré son efficacité, la Carte Vitale est contournée par les barrières financières résiduelles...}

\subsection*{\underline{2. Les Limites face aux Inégalités Sociales}}

\IdeeP{Les inégalités se manifestent par un taux élevé de renoncement aux soins, directement lié au revenu.}
\SousAxe{Le Reste à Charge (RAC) comme Barrière}
\begin{itemize}
    \item **Dépassements d'Honoraires :** Pratiqués en Secteur 2, le RAC est trop élevé sans une mutuelle coûteuse.
    \item **Soins Spécifiques :** Fort renoncement sur l'Optique et le Dentaire, malgré le dispositif "100\% Santé" (dont les équipements peuvent ne pas suffire).
    \item **Cumul des Désavantages :** Les personnes pauvres ont plus de pathologies, mais moins de ressources pour y faire face.
\end{itemize}
\ChiffreCle{Espérance de Vie : 13 ans d'écart chez les hommes entre les plus riches et les plus modestes (Insee).}

\Transition{Ces obstacles financiers se combinent aux difficultés d'ordre géographique et matériel...}


\subsection*{\underline{4. Les Déterminants Sociaux}}

\IdeeP{L'état de santé est d'abord déterminé par la place dans la hiérarchie sociale, bien avant l'accès aux soins.}
\SousAxe{Les Facteurs Structurels Profonds}
\begin{itemize}
    \item **Éducation :** Un niveau d'éducation faible est corrélé à une moins bonne espérance de vie (choix de vie, compréhension des risques, navigation dans le système).
    \item **Emploi/Conditions de Travail :** Exposition aux risques professionnels, stress social et usure précoce du corps (ouvriers).
    \item **Logement :** Le mal-logement (humidité, froid, insalubrité) est un facteur direct de maladie et de stress chronique.
    \item **Fracture Administrative :** Difficulté d'accès à la Carte Vitale et aux droits pour les personnes sans domicile ou les plus isolées (risque d'exclusion par le numérique).
\end{itemize}
\ChiffreCle{Espérance de vie sans incapacité : 10 ans d'écart entre cadre et ouvrier (34 ans vs 24 ans à 35 ans).}

% ----------------------------------------------------------------------
% FIN DU DOCUMENT
% ----------------------------------------------------------------------

\end{document}